\chapter{Administración del plan}
\label{cap:administracion}

\section{Cronograma}

La \Cref{tab:cronograma} resume las fases de elaboración del plan de tesis (Guía~N°01) y la
asignación de responsables. Las fases 0--E corresponden al avance documentado al cierre de mayo
2026; las fases F--H concentran la metodología, la administración y la revisión de cierre antes de
la entrega al docente y a la plataforma del curso.

\begin{table}[H]
  \centering
  \caption{Cronograma de elaboración del plan de tesis}
  \label{tab:cronograma}
  \small
  \begin{tabular}{@{}p{0.06\textwidth}p{0.34\textwidth}p{0.12\textwidth}p{0.12\textwidth}p{0.22\textwidth}@{}}
    \toprule
    \textbf{Fase} & \textbf{Actividad principal} & \textbf{Inicio} & \textbf{Fin} & \textbf{Responsable} \\
    \midrule
    0 & Bibliografía: fichas + Anexo B (14 papers) & Mar 2026 & Abr 2026 & Jack \\
    A & Esqueleto LaTeX e índices automáticos & Abr 2026 & Abr 2026 & Ambos \\
    B & Cap.~1 planteamiento + Fig.~1--3 & Abr 2026 & May 2026 & Alex \\
    C & Cap.~3 §3.5--§3.6 estado del arte & May 2026 & May 2026 & Alex \\
    D & Cap.~2 objetivos, hipótesis, variables & May 2026 & May 2026 & Ambos \\
    E & Cap.~3 §3.1--§3.4 marco + glosario & May 2026 & May 2026 & Jack \\
    F & Cap.~4 diseño metodológico + Anexo F & May 2026 & Jun 2026 & Jack \\
    G & Cap.~5 administración + anexos LaTeX & May 2026 & Jun 2026 & Jack \\
    H & Revisión de cierre y entrega docente & Jun 2026 & Jun 2026 & Ambos \\
    \bottomrule
  \end{tabular}
\end{table}


\section{Presupuesto}

El plan de tesis es un producto documental del curso de Proyecto de Investigación I; no requiere
inversión en infraestructura productiva. El presupuesto estimado es \textbf{cero soles} en licencias
y despliegue, dado el uso de herramientas libres (\LaTeX{}, Python, Zotero, Git).

\begin{table}[H]
  \centering
  \caption{Presupuesto estimado del plan de tesis}
  \label{tab:presupuesto}
  \small
  \begin{tabular}{@{}p{0.42\textwidth}rr@{}}
    \toprule
    \textbf{Ítem} & \textbf{Costo (S/)} & \textbf{Fuente} \\
    \midrule
    Licencias de software (\LaTeX{}, Python, Zotero, Git) & 0 & Open source \\
    Acceso a bases bibliográficas & 0 & Suscripción institucional \\
    Infraestructura cloud / despliegue API & 0 & Fuera de alcance §1.6 \\
    Impresión del documento (opcional) & $\leq$ 30 & Recursos personales \\
    \midrule
    \textbf{Total estimado} & \textbf{$\leq$ 30} & \\
    \bottomrule
  \end{tabular}
\end{table}

El costo principal del proyecto es el \textbf{tiempo de dedicación} de los integrantes
($\approx$ 8--10 horas semanales durante el ciclo académico), distribuido según el cronograma de la
\Cref{tab:cronograma}.

\section{Financiamiento}

La elaboración del plan de tesis no cuenta con financiamiento externo ni contratos de investigación.
Las actividades se sostienen con los siguientes medios:

\begin{itemize}
  \item \textbf{Recursos propios} de los integrantes (equipos personales, conectividad e impresión
        opcional del documento).
  \item \textbf{Recursos institucionales} de la Universidad Nacional de Ingeniería: acceso a bases
        bibliográficas (IEEE Xplore, ACM Digital Library y repositorios académicos) mediante la
        matrícula de posgrado; uso de aula virtual y sesiones de asesoría del curso Proyecto de
        Investigación I.
  \item \textbf{Software libre} (\LaTeX{}, Python, Zotero, Git) sin costo de licenciamiento,
        coherente con el presupuesto de la \Cref{tab:presupuesto}.
\end{itemize}

No se prevén viáticos, contratación de personal ni adquisición de hardware adicional para el alcance
delimitado en §1.6 (solo plan de tesis, sin experimentación a escala ni despliegue productivo).

\section{Recursos}

\subsection{Recursos humanos}

\begin{itemize}
  \item \textbf{Alex Mancilla:} planteamiento del problema, Cap.~1, revisión cruzada de coherencia
        narrativa y apoyo en exposiciones.
  \item \textbf{Jack Paitan:} sistematización bibliográfica (Fase~0), marco teórico §3.1--§3.4,
        diseño metodológico (Cap.~4), administración (Cap.~5), anexos y matrices Excel.
  \item \textbf{Docente del curso:} retroalimentación (sesiones PI-03 a PI-10) y validación del
        hilo conductor antes de la entrega formal.
\end{itemize}

\subsection{Recursos materiales y de software}

\begin{itemize}
  \item \textbf{Hardware:} equipos personales con acceso a internet (macOS); sin infraestructura
        de servidor propia.
  \item \textbf{Software:} \LaTeX{} (MacTeX), Python 3 + \texttt{openpyxl}/\texttt{matplotlib},
        editor de texto/Git para colaboración, Zotero para gestión bibliográfica.
  \item \textbf{Fuentes documentales:} contratos OpenAPI~3.x bancarios públicos o de referencia;
        especificaciones \ServiceDomain{} BIAN; repositorio \texttt{bibliografia/} del proyecto.
  \item \textbf{Acceso académico:} IEEE Xplore, ACM Digital Library y bases institucionales de la
        universidad para descarga de papers del estado del arte.
\end{itemize}

\subsection{Recursos bibliográficos}

Catorce papers de estado del arte (Anexo~B), veintidós papers de aporte para sustentar técnicas
(Cap.~4) y siete referencias de contexto OpenAPI--BIAN. La cadena de búsqueda y criterios de
selección se documentan en el Anexo~A.

\section{Roles del equipo}

\begin{table}[H]
  \centering
  \caption{Distribución de roles del equipo}
  \label{tab:roles}
  \small
  \begin{tabular}{@{}p{0.22\textwidth}p{0.68\textwidth}@{}}
    \toprule
    \textbf{Integrante} & \textbf{Responsabilidad principal} \\
    \midrule
    Alex Mancilla &
    Cap.~1 (§1.1--§1.9); coherencia problema--objeto--referencia BIAN; exposición oral;
    revisión cruzada Cap.~2--3 \\
    Jack Paitan &
    Fase~0; §3.5--§3.6 y §3.1--§3.4; Cap.~2--4; Cap.~5; anexos A--F; matrices Excel;
    compilación PDF y carga en aula virtual \\
    Ambos &
    Revisión de hilo conductor; matriz de consistencia; cierre Fase H; decisiones de alcance
    (Guía~N°01) \\
    \bottomrule
  \end{tabular}
\end{table}

\section{Riesgos}

\begin{table}[H]
  \centering
  \caption{Matriz de riesgos del plan de tesis}
  \label{tab:riesgos}
  \small
  \begin{tabular}{@{}p{0.22\textwidth}p{0.10\textwidth}p{0.10\textwidth}p{0.44\textwidth}@{}}
    \toprule
    \textbf{Riesgo} & \textbf{Prob.} & \textbf{Impacto} & \textbf{Mitigación} \\
    \midrule
    Ambigüedad en definición operativa de S o E &
    Media & Alta &
    Documentar fórmulas y umbrales en Cap.~4; revisión cruzada Alex/Jack; validación docente \\
    Retraso en redacción Cap.~4--5 &
    Media & Media &
    Cronograma \Cref{tab:cronograma}; priorizar diseño reproducible sin experimentación \\
    Inconsistencia entre matriz (Anexo C) y Cap.~2 &
    Baja & Alta &
    Regenerar Excel desde script; trazabilidad OG--OE en revisión Fase H \\
    Referencias incompletas o no uniformes &
    Media & Media &
    Export Zotero + lista APA en referencias; verificar citas Cap.~1--3 \\
    Sobrecarga de alcance (construcción/experimentación) &
    Baja & Alta &
    Delimitación §1.6; matriz sin piloto; sin OE de validación experimental \\
    \bottomrule
  \end{tabular}
\end{table}

